\usepackage[utf8]{inputenc}
\usepackage[T1]{fontenc}
\usepackage{graphicx}
\usepackage[x11names]{xcolor}
\usepackage{tikz}
\usepackage{fontawesome5}
\usepackage{bm}
\usepackage[most]{tcolorbox}
\usepackage{enumerate}
\usepackage[shortlabels]{enumitem}
\tcbuselibrary{skins,raster,theorems,breakable}

\usepackage{amsmath,amssymb,textcomp}
\usepackage{mathtools}
\everymath{\displaystyle}
\usepackage{multicol}
\usepackage{multirow}
\setlength{\columnseprule}{0pt}
\setlength{\columnsep}{20.0pt}


\usepackage{geometry}
\geometry{a4paper,left=10mm,right=10mm,top=10mm,bottom=15mm}

\newcommand{\trans}[1]{{#1}^{\mathsf{T}}}
\newcommand{\ev}[1]{\mathbb{E}\left[#1\right]}
\newcommand{\unmezz}{\frac{1}{2}}
	\newcommand{\Z}{\mathbb{Z}}
\newcommand{\R}{\mathbb{R}}
\newcommand{\N}{\mathbb{N}}
\newcommand{\T}{\mathbb{T}}
\newcommand{\Tbar}{\overline{\T}}
\newcommand{\Nstar}{\N^{*}}
\newcommand{\Rext}{\overline{\R}}
\newcommand{\io}{\text{ i.o.}}
\newcommand{\equalexpl}[1]{%
	\underset{\substack{\uparrow\\\mathrlap{\text{\vspace{-3cm}\hspace{-1em}#1}}}}{=}}
\newcommand{\dif}{\mathop{}\!\mathrm{d}}
\newcommand{\convas}{\xrightarrow[]{\text{a.s.}}}
\newcommand{\convpr}{\xrightarrow[]{\pr}}
\newcommand{\convd}{\xrightarrow[]{\text{d}}}
\newcommand{\convlp}{\xrightarrow[]{\lp}}
\newcommand{\convw}{\xrightarrow[]{\mathrm{weak}}}
\newcommand{\as}{\text{ a.s.}}
\newcommand{\asstnr}{\sim N(0,1)}
%\def\checkmark{\tikz\fill[scale=0.4](0,.35) -- (.25,0) -- (1,.7) -- (.25,.15) -- cycle;} 
\newcommand{\dx}{\dif x}
\newcommand{\dy}{\dif y}
\newcommand{\dt}{\dif t}
\newcommand{\du}{\dif u}
\newcommand{\ds}{\dif s}
\newcommand{\dw}{\dif\omega}
\newcommand{\dmu}{\dif\mu}
\newcommand{\dpr}{\dif\pr}
\newcommand{\E}{\mathscr{E}}
\newcommand{\B}{\mathscr{B}}
\newcommand{\F}{\mathcal{F}}
\newcommand{\G}{\mathscr{G}}
\newcommand{\HS}{\mathscr{H}}
\newcommand{\Zn}{\mathscr{Z}}
\newcommand{\D}{\mathscr{D}}
\newcommand{\Es}{\mathscr{S}}
\newcommand{\xbar}{\overline{X}}
\newcommand{\rbar}{\overline{\R}}
\newcommand{\ybar}{\overline{Y}}
\newcommand{\xxbar}{\overline{\xbar}}
\newcommand{\xtilde}{\widetilde{X}}
\newcommand{\ifonly}{\underline{if and only if}}
%\newcommand{\Lp}{L^{p}}

\newtcolorbox{riquadro}[1][]{enhanced, breakable,frame style={top color=Aquamarine3,bottom color=Aquamarine3},colback=black!80,coltext=Aquamarine3!50,title=#1,fonttitle=\bfseries,halign title=center,coltitle=black!80}
\makeatletter
\renewcommand*{\maketitle}{%
	\noindent
	\begin{minipage}{0.4\textwidth}
		\begin{tikzpicture}
			\node[rectangle,rounded corners=6pt,inner sep=10pt,fill=SteelBlue4,text width= 0.95\textwidth] {\color{white}\Huge \@title};
		\end{tikzpicture}
	\end{minipage}
	\hfill
	\begin{minipage}{0.55\textwidth}
		\begin{tikzpicture}
			\node[rectangle,rounded corners=3pt,inner sep=10pt,draw=Turquoise4,text width= 0.95\textwidth] {\LARGE \@author};
		\end{tikzpicture}
	\end{minipage}
}%
\makeatother
\newtcbtheorem[number within=section]{myproof}{Proof}{
	enhanced, sharp corners, breakable,
	colframe=LightSalmon1, coltitle=black, colbacktitle=gray!24,
	coltext=LightSalmon1!60,
	colback=black!80,
	rounded corners,
	fonttitle=\bfseries,
	separator sign={:},
	theorem hanging indent/.try=0pt,
}{ciao}

% custom section
\usepackage[explicit]{titlesec}
\newcommand*\sectionlabel{}
\newcommand{\pr}{\mathbb{P}}
\newcommand{\mbf}[1]{\mathbf{#1}}

\titleformat{\section}
{\gdef\sectionlabel{}
	\normalfont\rmfamily\Large\bfseries}
{\gdef\sectionlabel{\thesection\ }}{0pt}
{
	\noindent
	\begin{tikzpicture}
		\node[rectangle,rounded corners=3pt,inner sep=4pt,fill=Turquoise4,text width= 0.95\columnwidth] {\color{white}\sectionlabel#1};
	\end{tikzpicture}
}
\titlespacing*{\section}{0pt}{0pt}{-10pt}
\newcommand*\subsectionlabel{}
\titleformat{\subsection}
{\gdef\subsectionlabel{}
	\normalfont\rmfamily\Large\bfseries}
{\gdef\subsectionlabel{\thesubsection\ }}{0pt}
{
	\noindent
	\begin{tikzpicture}
		\node[rectangle,inner sep=4pt,fill=Turquoise4!60!black,rounded corners=3pt,text width= 0.9\columnwidth] {\color{white}\subsectionlabel#1};
	\end{tikzpicture}
}
\titlespacing*{\subsection}{0pt}{0pt}{0pt}

% custom footer
\usepackage{fancyhdr}
\makeatletter
\pagestyle{fancy}
\fancyhead{}
\fancyfoot[C]{\footnotesize \textcopyright\ \@date\ \ \@author}
\renewcommand{\headrulewidth}{0pt}
\renewcommand{\footrulewidth}{0pt}
\makeatother
